% §4 Main result.  Delivers one theorem: p = Theta(log m / m) suffices for the entry-wise
% criterion eq:entrywise to hold w.h.p.  The statement is about the criterion, not about
% which split of the tree is returned -- that is §5.
\section{Main Result}\label{sec:main}
The entry-wise recovery criterion \eqref{eq:entrywise} is what the sampling rate
has to buy. This section converts it into a sufficient condition on $p$ for the
Laplacian $L=D-S$, stated in the structural quantities of \Cref{sec:model}: the
margin $\rho$ of \eqref{eq:margin_definition} and the imbalance $\eta$ of
\eqref{eq:imbalance_ratio}.

\begin{theorem}[Exact recovery of the full-data split]\label{thm:main-sim}
Let $S$ satisfy the non-balanced CFBM of \Cref{def:cfbm} with $m$ terminal nodes and
imbalance $\eta\ge1$, let $\rho$
be the structural margin and let $\vvh$ be the empirical Fiedler vector --- the
unit eigenvector belonging to the second-smallest eigenvalue of the sampled
Laplacian $\hat L=\hat D-\hat S$, formed from entries observed with probability
$p$ as in the observation model \eqref{eq:sampling_estimator}. Assume
\Cref{asm:margin}. For the entry-wise recovery criterion
\eqref{eq:entrywise} to hold \hp, i.e.\
$\norm{\vvh-\vv}_\infty<\min_i|\vv_i|$, it suffices that
\begin{equation}\label{eq:main_rate}
   p\;>\;\frac{8(1+\eta)^3\,\bze^2\,\log m}{m\,(\rho-\eta\,\Sout^{\max})^2}
       \;=\;\Theta\!\left(\frac{\log m}{m}\right).
\end{equation}
\end{theorem}

\noindent \eqref{eq:main_rate} is a \emph{sufficient} condition. No matching lower
bound is proved anywhere in this paper, so the $\Theta(\log m/m)$ is the growth
order of that sufficient bound and not a claim that the rate is tight or
order-optimal.

\noindent The distance operator $\Bop=H\Dmat H$ of \Cref{sec:problem} admits a
statement of the same shape --- entry-wise sign agreement of an empirical
structural eigenvector against a piecewise-constant target, at a rate governed by
a spectral gap. Since the hypotheses of the two operators are not nested, it is a
\emph{sibling} of \Cref{thm:main-sim} rather than a corollary of it;
\Cref{app:dist} states it as \Cref{thm:main-dist} and proves it.
